\section{Zaključak}

Razvoj informacionog sistema pozorišta predstavlja značajan korak ka modernizaciji poslovanja jedne ovakve kulturne institucije. Kroz detaljnu analizu zahteva, modelovanje slučajeva upotrebe, podataka i poslovnih procesa, projektovano je sveobuhvatno rešenje koje pokriva sve ključne aspekte rada pozorišta – od kreiranja predstava i upravljanja repertoarom, preko prodaje karata, do finansijske administracije i upravljanja ljudskim resursima.

Implementacija ovakvog sistema omogućava višestruke benefite:
\begin{itemize}
    \item \textbf{Povećanje efikasnosti}: Automatizacija rutinskih zadataka, poput prodaje karata i generisanja finansijskih izveštaja.
    \item \textbf{Transparentnost i kontrola}: Centralizovana baza podataka o finansijama, zaposlenima i tehničkim resursima omogućava bolji uvid u poslovanje i olakšava donošenje odluka.
    \item \textbf{Poboljšano korisničko iskustvo}: Online sistem za kupovinu karata i lako dostupne informacije o repertoaru značajno unapređuju iskustvo publike i čine pozorište dostupnijim.
\end{itemize}

Iako trenutno rešenje pokriva primarne potrebe pozorišta, arhitektura sistema je dizajnirana tako da bude modularna i otvorena za buduće nadogradnje. Među potencijalnim proširenjima ističe se razvoj mobilne aplikacije koja bi dodatno olakšala interakciju sa publikom.

Verujemo da ovaj sistem uspešno modernizuje administrativni deo posla, čime se zaposlenima otvara prostor da se više posvete onome što čini suštinu pozorišta – umetnosti i publici.
